% This document is compiled using pdfLaTeX
% You can switch XeLaTeX/pdfLaTeX/LaTeX/LuaLaTeX in Settings

\documentclass{article}
\usepackage{ctex}
\usepackage{amsmath} % 确保引入了此宏包以支持 aligned

\title{因子日历2026}

\date{\today}

\begin{document}

\maketitle

\section{超大单买入占比(高频因子-资金流类)}

\subsection{因子定义}

\begin{equation}
    \text{超大单买入占比}_{i,t} = \frac{\mathrm{SuperBigBuy}_{i,t}}{\mathrm{SuperBigBuy}_{i,t} + \mathrm{SuperBigSell}_{i,t}}
\end{equation}

\subsection{变量定义}
\begin{itemize}
    \item ① 按照委买单号和委卖单号对 level2 逐笔成交数据中的成交金额（已成交部分）进行汇总求和得到当天股票 \( i \) 所有买单和卖单的成交金额；
    \item ② 对股票 \( i \)，将订单成交金额占当天总成交额比例超过 1\% 的订单定义为超大单（阈值在 0.5\%~2\% 区间较好）；
    \item ③ \( \mathrm{SuperBigBuy}_{i,t} \) 和 \( \mathrm{SuperBigSell}_{i,t} \) 分别为股票 \( i \) 在交易日 \( t \) 的所有超大买单、超大卖单的成交金额；
    \item ④ 因为超大单买卖金额在大票中不稳定，所以在计算全天、早盘超大单买入占比因子时，分子、分母均用 20 天超大单金额的均值替代单日的超大单金额。
\end{itemize}

\subsection{说明}
超大单刻画了潜在的过度投机和流动性冲击的负向 Alpha 影响，通常与未来收益负相关，超大买单占比越大，未来反转几率越大，小市值的大单反转更明显。

\subsection{参考文献}
\begin{enumerate}
    \item 朱剑涛, 王星星. 超大单冲击对大单因子的影响, 2022, 东方证券.
\end{enumerate}

\section{基于涨跌停的注意力捕捉(行为金融因子-投资者注意力)}

\subsection{因子定义}

将涨跌停股票的数量除以当日交易的股票数得到每日全市场涨跌停的比率作为市场关注代理指标：
\begin{equation}
    \mathrm{LIMIT\_PROP}_d = \frac{n_{\mathrm{upper},d} + n_{\mathrm{lower},d}}{N_d}
\end{equation}

将市场关注代理指标与股票日收益进行月度时序回归，代理指标新的系数绝对值 \(|\beta_{i,t}|\) 即为股票对市场关注的敏感程度：
\begin{equation}
    r_{i,d} = \mu_{i,t} + \beta_{i,t} \mathrm{LIMIT\_PROP}_d + \rho_{1,i,t} \mathrm{Mkt}_d + \rho_{2,i,t} \mathrm{SMB}_d + \rho_{3,i,t} \mathrm{HML}_d + \rho_{4,i,t} \mathrm{UMD}_d + \varepsilon_{i,t}
\end{equation}

\subsection{变量定义}
\begin{itemize}
    \item ① \(n_{\mathrm{upper},d}\)、\(n_{\mathrm{lower},d}\)、\(N_d\) 分别为 \(d\) 日涨停股票数量、跌停股票数量、当日交易股票总数；
    \item ② \(r_{i,d}\) 为个股日度收益；
    \item ③ \(\mathrm{Mkt}_d\)、\(\mathrm{SMB}_d\)、\(\mathrm{HML}_d\) 为经典的 Fama-French 三因子；
    \item ④ \(\mathrm{UMD}_d\) 为动量因子收益，即全市场过去 11 个月累积收益前后 30\% 的组合在 \(d\) 日的收益差。
\end{itemize}

\subsection{说明}
不同股票对同一市场关注度事件的反映程度往往是不同的，可借助个股收益对市场关注度事件的敏感程度来衡量股票对该事件的注意力捕捉效应；捕捉能力越强（敏感程度越高），表明个股更能吸引市场注意，从而导致买盘压力增大，股价高估；上述因子以涨跌停作为市场关注度事件，占比越高说明市场关注度越高。

\subsection{参考文献}
\begin{enumerate}
    \item Lin, F., Qiu, Z., \& Zheng, W. (2023). \textit{Cranes among chickens: The general-attention-grabbing effect of daily price limits in China's stock market}. Journal of Banking \& Finance, 150, 106818.
    \item 陈丹妮. 2024. 投资者有限关注及注意力捕捉与溢出. 中信出版社.
\end{enumerate}

\section{尾盘成交占比(高频因子-成交分布类)}

\subsection{因子定义}
\begin{equation}
    \frac{1}{T} \sum_{n=t}^{t-T+1} \frac{\mathrm{Vol}_{i,14:30-15:00,n}}{\mathrm{Vol}_{i,n}}
\end{equation}

\subsection{变量定义}
\begin{itemize}
    \item ① \( \mathrm{Vol}_{i,14:30-15:00,n} \) 为第 \( i \) 只股票在第 \( n \) 个交易日内 14:30-15:00 尾盘的成交量；
    \item ② \( \mathrm{Vol}_{i,n} \) 为第 \( i \) 只股票在第 \( n \) 个交易日的总成交量；
    \item ③ 月度选股下，\( T=20 \) 个交易日；周度选股下，\( T=5 \) 个交易日。
\end{itemize}

\subsection{说明}
尾盘成交量占比因子和股票未来收益负相关，其较好的选股效果可能源于：1. 尾盘投机度高，容易出现价格操纵行为；2. 非知情交易者（散户）不愿承担日内波动，更倾向于尾盘交易，而知情交易者（机构）则倾向于在早盘交易。

\subsection{参考文献}
\begin{enumerate}
    \item 冯佳睿等. 2020. 高频因子的现实与幻想. 海通证券.
\end{enumerate}

\section{订单失衡(高频因子-流动性类)}

\subsection{因子定义}
\begin{equation}
    \mathrm{VOI}_t = \Delta V_t^{\mathrm{WB}} - \Delta V_t^{\mathrm{WA}}
\end{equation}
其中，
\begin{equation}
    \Delta V_t^{\mathrm{WB}} = 
    \begin{cases}
        V_t^{\mathrm{WB}} , & P_t^{\mathrm{B}} > P_{t-1}^{\mathrm{B}} \\
        V_t^{\mathrm{WB}} - V_{t-1}^{\mathrm{WB}}, & P_t^{\mathrm{B}} = P_{t-1}^{\mathrm{B}} \\
        0, & P_t^{\mathrm{B}} < P_{t-1}^{\mathrm{B}}
    \end{cases}
\end{equation}
\begin{equation}
    \Delta V_t^{\mathrm{WA}} = 
    \begin{cases}
        V_t^{\mathrm{WA}}, & P_t^{\mathrm{A}} < P_{t-1}^{\mathrm{A}} \\
        V_t^{\mathrm{WA}} - V_{t-1}^{\mathrm{WA}}, & P_t^{\mathrm{A}} = P_{t-1}^{\mathrm{A}} \\
        0, & P_t^{\mathrm{A}} > P_{t-1}^{\mathrm{A}}
    \end{cases}
\end{equation}

\subsection{变量定义}
\begin{itemize}
    \item ① \( V^{\mathrm{WB}(\mathrm{WA})} = \sum w_i \times V_{i,t}^{\mathrm{B}(\mathrm{A})} \)，\( w_i = 1 - (i-1)/5, i=1,2,3,4,5 \) 为衰减加权委托量；
    \item ② \( P_t^{\mathrm{B}} \) 和 \( P_t^{\mathrm{A}} \) 分别为 \( t \) 时刻的买一价和卖一价；
    \item ③ \( V_t^{\mathrm{B}} \) 和 \( V_t^{\mathrm{A}} \) 分别为 \( t \) 时刻的买一和卖一的委托量。
\end{itemize}

\subsection{说明}
订单失衡因子为委托买入量与委托卖出量之差，衡量了盘口订单的不平衡程度，刻画了交易者的交易意图（知情交易者的信息差异使得订单不平衡）；订单失衡因子短期与收益率呈正相关，长期会出现反转（存在散户的追涨杀跌和主力的短期市场操纵）。

\subsection{参考文献}
\begin{enumerate}
    \item 丁鲁明, 陈升锐. 2020. 高频量价选股因子初探. 中信建投证券.
\end{enumerate}

\section{CSAD 模型（基于“小型”板块）(行为金融因子-羊群效应)}

\subsection{因子定义}

1. 每个月末提取当前可交易股票过去 120 个交易日的涨跌幅数据，计算股票 \( S_t \) 和其他股票（不活跃交易日不超过 20 天）之间日涨跌幅相关系数；

2. 选取相关系数最大的 9 只股票，连同 \( S_t \) 一起构成一个板块：
   \begin{equation}
       S_i \ (i = 1, \dots, 10)
   \end{equation}

3. 计算板块过去 120 个交易日的横截面绝对偏差 CSAD：
   \begin{equation}
       \mathrm{CSAD}_t = \frac{1}{10} \sum_{i=1}^{10} |r_{it} - \bar{r}_t|
   \end{equation}

4. 对 CSAD 进行时序标准化，并计算过去 20 个交易日的均值作为羊群效应因子：
   \begin{equation}
       Z_{jt} = \frac{\mathrm{CSAD}_{jt} - \mu_{\mathrm{CSAD}_j}}{\sigma_{\mathrm{CSAD}_j}}
   \end{equation}
   \begin{equation}
       f_{jt} = -\frac{1}{20} \sum_{k=1}^{20} Z_{j,t-k+1}
   \end{equation}

\subsection{说明}
CSAD 利用股票收益在横截面的分散度来衡量市场羊群效应的强弱；为了让因子对个股有选股能力，将与个股走势相近的股票组成“小型板块”，借助小板块上的羊群效应的强弱程度来近似反映该股票近期关注度的变化。因子值越大，板块近期的 CSAD 越小（即板块羊群效应越强），所在板块受关注程度提升越明显，未来收益越高。

\subsection{参考文献}
\begin{enumerate}
    \item Chang, E. C., Cheng, J., \& Khorana, A. (2000). \textit{An Examination of Herd Behavior in Equity Markets: An International Perspective}. Journal of Banking \& Finance, 24(10), 1651-1679.
    \item 丁鲁明, 陈元鹏. 2019. 基于市场羊群效应的股票 alpha 研究. 中信出版社.
\end{enumerate}

\section{耀眼波动率(高频因子-波动跳跃类)}

\subsection{因子定义}

1. 对股票 \( i \)，计算 t 日内 1 分钟成交量序列的环比差值，将成交量差值大于“该日差值序列均值+1倍标准差”的时刻定义为“成交量激增时刻 n”；

2. 利用分钟收盘价，计算股票 \( i \) 在 t 日内 1 分钟收益率序列；计算成交量激增时刻及其随后 4 分钟 \([n, n+4]\) 区间内 1 分钟收益率的标准差，即为“耀眼波动率”；

3. 将 t 日内所有耀眼波动率求均值即为“日耀眼波动率”；

4. 将各股票的“日耀眼波动率”减去该指标的横截面均值再取绝对值，即为“适度日耀眼波动率”。

\subsection{变量定义}
\begin{itemize}
    \item ① 计算因子需剔除开盘和收盘数据，仅考虑日内分钟频数据；
    \item ② 可对最近 20 个交易日的日度“适度日耀眼波动率”求均值和标准差，并等权合成，即可得“月耀眼波动率”因子。
\end{itemize}

\subsection{说明}
“耀眼波动率”衡量了日内成交量激增所引起的价格波动程度；“耀眼波动率”与未来收益负相关，若成交量激增后的价格波动较小，说明利好消息未被市场广泛认可，适度承担该风险能获得风险补偿；若引起的价格波动很大，可能代表投资者已反应过度，未来会补跌。

\subsection{参考文献}
\begin{enumerate}
    \item 曹春晓. 2022. 成交额激增时刻蕴含的 alpha 信息. 方正证券.
\end{enumerate}

\section{每笔成交量收益率相关性(高频因子-量价相关性类)}

\subsection{因子定义}
\begin{equation}
     \frac{\sum w_i \times ret_i}{\mathrm{std}(\mathrm{pvol}_i) \times \mathrm{std}(\mathrm{ret}_i)} = \mathrm{corr}(\mathrm{pvol}, \mathrm{ret})
\end{equation}

\begin{equation}
     w_i = \mathrm{pvol}_i - mean(\mathrm{pvol_i})
\end{equation}

\subsection{变量定义}
\begin{itemize}
    \item ① \(\mathrm{ret}_i\) 和 \(\mathrm{pvol}_i\) 分别为每个时间段的收益率和每笔成交量=成交量/成交笔数，时间段为分钟频率。
\end{itemize}

\subsection{说明}
每笔成交量收益率相关性本质为以每笔成交量与自身均值之差为权重对收益率序列进行加权，是利用每笔成交量对反转因子做了“提纯”，相关系数相比局部筛选保留了全时间段的所有信息。每笔成交量较大、呈现反转效应的时间段收益率的权重大于0，且成交量越大，权重越大，进而加强反转效应；反之，反向加强动量效应。

\subsection{参考文献}
\begin{enumerate}
    \item 覃川桃, 郑起. 2021. 高频因子（十）：量价关系中的反转微观结果. 长江证券.
\end{enumerate}

\section{海外运营信息(图谱网络-动量溢出)}

\subsection{因子定义}
\begin{equation}
    \mathrm{InfoProxy}_{i,j,t-1}{\mathrm{(Foreign)}} = \sum_{c\nep U.S.} f_{i,t-1}^c R_{j,t-1}^c
\end{equation}

\subsection{变量定义}
\begin{itemize}
    \item ① \( f_{i,t-1}^c \) 为跨国公司在 \( c \) 国家的销售额占比；
    \item ② \( R_{j,t-1}^c \) 为 \( c \) 国家的 \( j \) 行业回报（从 Datastream 全球股权部门指数中获取跨国公司运营所在国家的行业回报数据）。
\end{itemize}

\subsection{说明}
海外运营信息因子反映了跨国公司海外业务活动情况，跨国公司海外业务所对应的各国行业回报信息可以用来预测跨国公司自身股价的未来收益（因子方向为正）；投资者对海外运营信息的关注度较低，存在反应不足，进而导致股价未能及时反映跨国公司的实际经营状况变化，投资者可通过捕捉这些股价尚未完全反映的新信息而获得超额收益。

\subsection{参考文献}
\begin{enumerate}
    \item Huang, Xing. (2015). \textit{Thinking outside the borders: investors' underreaction to foreign operations information}. Review of Financial Studies, 28(11), 3109-3152.
\end{enumerate}

\section{开盘后买入意愿占比(高频因子-资金流类)}

\subsection{因子定义}


\begin{equation}
    \text{开盘后买入意愿占比} = \frac{1}{T} \sum_{n=t}^{t-T+1} \frac{\sum_{j=1}^{N} \text{买入意愿}_{i,j,n}}{\sum_{j=1}^{N} \text{成交额}_{i,j,n}}
\end{equation}
\begin{equation}
    \text{买入意愿}_{i,j,n} = \text{净主买成交额}_{i,j,n} - \text{净委买增额}_{i,j,n}
\end{equation}
\begin{equation}
    \text{净主买成交额}_{i,j,n} = \text{主动买入成交额}_{i,j,n} - \text{主动卖出成交额}_{i,j,n}
\end{equation}
\begin{equation}
    \text{净委买增额}_{i,j,n} = \text{委托买单增加量}_{i,j,n} - \text{委托卖单增加量}_{i,j,n}
\end{equation}

\subsection{变量定义}
\begin{itemize}
    \item ① 净主买成交额由逐笔成交数据计算得到，具体可参考开盘后净主买占比/强度等因子；
    \item ② 净委买增额由盘口委托快照数据计算得到，具体可参考开盘后净委买增额占比因子；
    \item ③ \( i, j, n \) 分别表示第 \( i \) 只股票在第 \( n \) 个交易日内的第 \( j \) 分钟的数据；
    \item ④ 开盘后的买入意愿占比用的是 9:30-10:00 范围内的数据；
    \item ⑤ 月度选股下，\( T=20 \) 个交易日；周度选股下，\( T=5 \) 个交易日。
\end{itemize}

\subsection{说明}
净委买变化额刻画了投资者还未释放的买入意愿，而净主买成交额则刻画了投资者已经释放的买入意愿，两者结合则得到广义的投资者主动买入意愿；开盘后 30 分钟内买入意愿占比越高，投资者的买入意愿越强。

\subsection{参考文献}
\begin{enumerate}
    \item 冯佳睿, 林青. 2020. 基于直观逻辑和机器学习的高频数据低频化应用. 海通证券.
\end{enumerate}

\section{P型成交量分布}

\subsection{因子定义}
高频因子-成交分布类

1. 计算同价成交量：将日内相同分钟收盘价的成交量累加至一起，得到在当前价格上的成交量总和；

2. 将成交量累计最大的价格定义为该股当日的成交量支撑点 \( \mathrm{VolumeSupportPrice} \)，包含该支撑点的附近区域定义为成交量支撑区域 \( \mathrm{VolumeSupportArea} \)；

3. 逐步计算 \( \mathrm{VSP} \) 周围的成交量累计值（价格由近及远，成交量由大及小），该累计值与全天成交量总和的比值超 50\% 的最小区域，即为成交量支撑区域 \( \mathrm{VSA} \)，该区域上限价格为 \( \mathrm{VSA\_High} \)，下限价格为 \( \mathrm{VSA\_Low} \)；

4. 计算下限价格 \( \mathrm{VSA\_Low} \) 与当日最高价之间的差异即得到 \( \mathrm{vsa\_low2max} \)：


\subsection{说明}
若下限价格与当日最高价越接近，说明成交量支撑区域越接近日内的高价区域，则该股日内的成交量分布越接近于 P 型成交量分布，预期上涨；当个股价格出现急剧上涨然后盘整，或短期急剧下跌但日内整体处于高位时，通常会出现 P 型成交量曲线。

\subsection{参考文献}
\begin{enumerate}
    \item 郑兆磊. 2022. 成交量分布中的 Alpha. 兴业证券.
\end{enumerate}
\end{document}